\documentclass{article}
\usepackage{iftex}
\ifXeTeX
  \usepackage{fontspec}
  \usepackage{xeCJK}
  \setmonofont[Path=/Users/yupyom/.src2tex/fonts/hackgen/, BoldFont=HackGen-Bold.ttf]{HackGen-Regular.ttf}
  \setCJKmonofont[Path=/Users/yupyom/.src2tex/fonts/hackgen/, BoldFont=HackGen-Bold.ttf]{HackGen-Regular.ttf}
  \setCJKmainfont[Path=/usr/local/texlive/2026/texmf-dist/fonts/opentype/public/haranoaji/, Extension=.otf, BoldFont=HaranoAjiMincho-Bold.otf]{HaranoAjiMincho-Regular}
  \XeTeXlinebreaklocale ""
  \xeCJKsetup{CJKglue={},CJKecglue={}}
\fi
\ifLuaTeX
  \usepackage{luatexja-fontspec}
  \setmonofont[Path=/Users/yupyom/.src2tex/fonts/hackgen/, BoldFont=HackGen-Bold.ttf]{HackGen-Regular.ttf}
  \setmainjfont{HaranoAjiMincho-Regular.otf}
  \ltjsetparameter{kanjiskip=0pt, xkanjiskip=0pt}
\fi
\usepackage[a4paper, margin=2cm]{geometry}
\usepackage{graphicx}
\newdimen\charwd
{\tt\global\setbox0=\hbox{x}\global\charwd=\wd0}
\frenchspacing
\usepackage{fancyhdr}
\pagestyle{fancy}
\renewcommand{\headrulewidth}{0pt}
\fancyhf{}
\fancyhead[R]{\rm src2tex-go version 2.23}
% plain TeX compatibility
\makeatletter
\providecommand{\eqalign}[1]{\vcenter{\openup1\jot\m@th
  \ialign{\strut\hfil$\displaystyle{##}$&$\displaystyle{{}##}$\hfil\crcr#1\crcr}}}
\makeatother
\fancyfoot[R]{\rm\hfill samples{\tt /}hanoi.c\qquad page \thepage}
\begin{document}

\ifx\sevenrm\undefined
  \font\sevenrm=cmr7 scaled \magstep0
\fi

\newread\MyStyle
\openin\MyStyle=src2latex.s2t
\ifeof\MyStyle
  \closein\MyStyle
\else
  \input src2latex.s2t
  \closein\MyStyle
\fi

\def\mc{\relax}
\def\gt{\relax}
\def\sc{\scshape}

\tt\mc 

\noindent
\rm\mc {\tt /}{\tt *}\  \hrulefill \rm\mc \ {\tt *}

\noindent
\mbox{}\hfill

\noindent
\mbox{}\hfill

\noindent
\mbox{} {}% beginning of TeX mode

\centerline{\bf Towers of Hanoi}

\begin{quote}
\noindent
This program gives an answer to the following famous problem (towers of
Hanoi).
There is a legend that when one of the temples in Hanoi was constructed,
three poles were erected and a tower consisting of 64 golden discs was
arranged on one pole, their sizes decreasing regularly from bottom to top.
The monks were to move the tower of discs to the opposite pole, moving
only one at a time, and never putting any size disc above a smaller one.
The job was to be done in the minimum numbers of moves. What strategy for
moving discs will accomplish this optimum transfer?
\end{quote}

% end of TeX mode  \rm\mc 

\noindent
\mbox{}\hfill

\noindent
\mbox{}{\tt *}\  \hrulefill \rm\mc \ {\tt *}{\tt /}
\tt\mc 

\noindent
\mbox{}\hfill

\noindent
\mbox{}\rm\mc {\tt /}{\tt *}\  \hrulefill\ hanoi.c\ \hrulefill \rm\mc \ {\tt *}{\tt /}
\tt\mc 

\noindent
\mbox{}\hfill

\noindent
\mbox{}\hfill

\noindent
\mbox{}{\tt\#}include{\tt\mc \ }{\tt <}stdio.h{\tt >}

\noindent
\mbox{}{\tt\#}define{\tt\mc \ }ARRAY{\tt\mc \ }8{\tt\mc \ }{\tt\mc \ \ \ \ \ \ \ \ }{\tt\mc \ \ \ \ \ \ \ \ }{\tt\mc \ \ \ \ \ \ \ \ }\rm\mc {\tt /}{\tt *}\  {}disc の数 \hfill \rm\mc \ {\tt *}{\tt /}
\tt\mc 

\noindent
\mbox{}\hfill

\noindent
\mbox{}int{\tt\mc \ }disc[3][ARRAY];{\tt\mc \ \ \ \ \ }{\tt\mc \ \ \ \ \ \ \ \ }{\tt\mc \ \ \ \ \ \ \ \ }\rm\mc {\tt /}{\tt *}\  {}disc に関するデータの置き場所
					   \hfill \rm\mc \ {\tt *}{\tt /}
\tt\mc 

\noindent
\mbox{}\hfill

\noindent
\mbox{}void{\tt\mc \ }init{\tt\_}array(void){\tt\mc \ \ \ }{\tt\mc \ \ \ \ \ \ \ \ }{\tt\mc \ \ \ \ \ \ \ \ }\rm\mc {\tt /}{\tt *}\  {}disc に関するデータの初期化
					   \hfill \rm\mc \ {\tt *}{\tt /}
\tt\mc 

\noindent
\mbox{}{\tt\char'173}

\noindent
\mbox{}{\tt\mc \ \ \ \ }int{\tt\mc \ }j;

\noindent
\mbox{}\hfill

\noindent
\mbox{}{\tt\mc \ \ \ \ }for{\tt\mc \ }(j{\tt\mc \ }={\tt\mc \ }0;{\tt\mc \ }j{\tt\mc \ }{\tt <}{\tt\mc \ }ARRAY;{\tt\mc \ }++j){\tt\mc \ }{\tt\mc \ \ \ \ \ \ \ \ }\rm\mc {\tt /}{\tt *}\  {}[問題] この {\bf for} ループ
					   を抜け出した後、２次元配列
					   \hfill \rm\mc \ {\tt *}{\tt /}
\tt\mc 

\noindent
\mbox{}{\tt\mc \ \ \ \ \ \ \ \ }{\tt\mc \ \ \ \ \ \ \ \ }{\tt\mc \ \ \ \ \ \ \ \ }{\tt\mc \ \ \ \ \ \ \ \ }{\tt\mc \ \ \ \ \ \ \ \ }\rm\mc {\tt /}{\tt *}\  \null{\tt disc[][]} にはどのよう
					   なデータが代入されることに
					   \hfill \rm\mc \ {\tt *}{\tt /}
\tt\mc 

\noindent
\mbox{}{\tt\mc \ \ \ \ \ \ \ \ }{\tt\mc \ \ \ \ \ \ \ \ }{\tt\mc \ \ \ \ \ \ \ \ }{\tt\mc \ \ \ \ \ \ \ \ }{\tt\mc \ \ \ \ \ \ \ \ }\rm\mc {\tt /}{\tt *}\  {}なるか、具体的に述べよ。
					   \hfill \rm\mc \ {\tt *}{\tt /}
\tt\mc 

\noindent
\mbox{}{\tt\mc \ \ \ \ \ \ }{\tt\char'173}

\noindent
\mbox{}{\tt\mc \ \ \ \ \ \ \ \ }{\tt\mc \ \ }disc[0][j]{\tt\mc \ }={\tt\mc \ }ARRAY{\tt\mc \ }{\tt -}{\tt\mc \ }j;

\noindent
\mbox{}{\tt\mc \ \ \ \ \ \ \ \ }{\tt\mc \ \ }disc[1][j]{\tt\mc \ }={\tt\mc \ }0;

\noindent
\mbox{}{\tt\mc \ \ \ \ \ \ \ \ }{\tt\mc \ \ }disc[2][j]{\tt\mc \ }={\tt\mc \ }0;

\noindent
\mbox{}{\tt\mc \ \ \ \ \ \ }{\tt\char'175}

\noindent
\mbox{}{\tt\char'175}

\noindent
\mbox{}\hfill

\noindent
\mbox{}void{\tt\mc \ }print{\tt\_}result(void){\tt\mc \ }{\tt\mc \ \ \ \ \ \ \ \ }{\tt\mc \ \ \ \ \ \ \ \ }\rm\mc {\tt /}{\tt *}\  {}結果の表示 \hfill \rm\mc \ {\tt *}{\tt /}
\tt\mc 

\noindent
\mbox{}{\tt\char'173}

\noindent
\mbox{}{\tt\mc \ \ \ \ }static{\tt\mc \ }long{\tt\mc \ }counter{\tt\mc \ }={\tt\mc \ }0;{\tt\mc \ \ \ \ }{\tt\mc \ \ \ \ \ \ \ \ }\rm\mc {\tt /}{\tt *}\  {}[問題] ここで {\tt static}
					   と宣言したのは、なぜなのか
					   \hfill \rm\mc \ {\tt *}{\tt /}
\tt\mc 

\noindent
\mbox{}{\tt\mc \ \ \ \ \ \ \ \ }{\tt\mc \ \ \ \ \ \ \ \ }{\tt\mc \ \ \ \ \ \ \ \ }{\tt\mc \ \ \ \ \ \ \ \ }{\tt\mc \ \ \ \ \ \ \ \ }\rm\mc {\tt /}{\tt *}\  {}考えよ。もしも {\tt static}
					   を削除したら、何が起こる \hfill \rm\mc \ {\tt *}{\tt /}
\tt\mc 

\noindent
\mbox{}{\tt\mc \ \ \ \ \ \ \ \ }{\tt\mc \ \ \ \ \ \ \ \ }{\tt\mc \ \ \ \ \ \ \ \ }{\tt\mc \ \ \ \ \ \ \ \ }{\tt\mc \ \ \ \ \ \ \ \ }\rm\mc {\tt /}{\tt *}\  {}か、実験をしてみよう。
					   \hfill \rm\mc \ {\tt *}{\tt /}
\tt\mc 

\noindent
\mbox{}{\tt\mc \ \ \ \ }int{\tt\mc \ }i,{\tt\mc \ }j;

\noindent
\mbox{}\hfill

\noindent
\mbox{}{\tt\mc \ \ \ \ }printf({\tt "}{\tt -}{\tt -}{\tt -}{\tt\#}{\tt\%}ld{\tt -}{\tt -}{\tt -}{\tt\char92}n{\tt "},{\tt\mc \ }++counter);

\noindent
\mbox{}{\tt\mc \ \ \ \ }for{\tt\mc \ }(i{\tt\mc \ }={\tt\mc \ }0;{\tt\mc \ }i{\tt\mc \ }{\tt <}={\tt\mc \ }2;{\tt\mc \ }++i)

\noindent
\mbox{}{\tt\mc \ \ \ \ \ \ }{\tt\char'173}

\noindent
\mbox{}{\tt\mc \ \ \ \ \ \ \ \ }{\tt\mc \ \ }printf({\tt "}[{\tt\%}d]{\tt\mc \ }{\tt "},{\tt\mc \ }i);

\noindent
\mbox{}{\tt\mc \ \ \ \ \ \ \ \ }{\tt\mc \ \ }for{\tt\mc \ }(j{\tt\mc \ }={\tt\mc \ }0;{\tt\mc \ }j{\tt\mc \ }{\tt <}{\tt\mc \ }ARRAY;{\tt\mc \ }++j)

\noindent
\mbox{}{\tt\mc \ \ \ \ \ \ \ \ }{\tt\mc \ \ \ \ }{\tt\char'173}

\noindent
\mbox{}{\tt\mc \ \ \ \ \ \ \ \ }{\tt\mc \ \ \ \ \ \ \ \ }if{\tt\mc \ }(disc[i][j]{\tt\mc \ }!={\tt\mc \ }0)

\noindent
\mbox{}{\tt\mc \ \ \ \ \ \ \ \ }{\tt\mc \ \ \ \ \ \ \ \ }{\tt\mc \ \ \ \ }printf({\tt "}{\tt\%}d{\tt\mc \ }{\tt "},{\tt\mc \ }disc[i][j]);

\noindent
\mbox{}{\tt\mc \ \ \ \ \ \ \ \ }{\tt\mc \ \ \ \ \ \ \ \ }else

\noindent
\mbox{}{\tt\mc \ \ \ \ \ \ \ \ }{\tt\mc \ \ \ \ \ \ \ \ }{\tt\mc \ \ \ \ }break;

\noindent
\mbox{}{\tt\mc \ \ \ \ \ \ \ \ }{\tt\mc \ \ \ \ }{\tt\char'175}

\noindent
\mbox{}{\tt\mc \ \ \ \ \ \ \ \ }{\tt\mc \ \ }printf({\tt "}{\tt\char92}n{\tt "});

\noindent
\mbox{}{\tt\mc \ \ \ \ \ \ }{\tt\char'175}

\noindent
\mbox{}{\tt\char'175}

\noindent
\mbox{}\hfill

\noindent
\mbox{}static{\tt\mc \ }int{\tt\mc \ }{\tt *}ptr[3];{\tt\mc \ \ \ \ \ }{\tt\mc \ \ \ \ \ \ \ \ }{\tt\mc \ \ \ \ \ \ \ \ }\rm\mc {\tt /}{\tt *}\  {}disc 移動用ポインタ \hfill \rm\mc \ {\tt *}{\tt /}
\tt\mc 

\noindent
\mbox{}\hfill

\noindent
\mbox{}void{\tt\mc \ }move{\tt\_}one{\tt\_}disc(int{\tt\mc \ }i,{\tt\mc \ }int{\tt\mc \ }j){\tt\mc \ \ \ \ \ \ \ \ }\rm\mc {\tt /}{\tt *}\  {}1枚の disc を pole $i$ から
					   pole $j$ に移動する \hfill \rm\mc \ {\tt *}{\tt /}
\tt\mc 

\noindent
\mbox{}{\tt\char'173}

\noindent
\mbox{}{\tt\mc \ \ \ \ }({\tt *}ptr[j]++){\tt\mc \ }={\tt\mc \ }({\tt *}{\tt -}{\tt -}ptr[i]);{\tt\mc \ \ }{\tt\mc \ \ \ \ \ \ \ \ }\rm\mc {\tt /}{\tt *}\  {}[問題] {\tt ++} はポインタ
					   の後にあり、{\tt --} はポインタ
					   \hfill \rm\mc \ {\tt *}{\tt /}
\tt\mc 

\noindent
\mbox{}{\tt\mc \ \ \ \ \ \ \ \ }{\tt\mc \ \ \ \ \ \ \ \ }{\tt\mc \ \ \ \ \ \ \ \ }{\tt\mc \ \ \ \ \ \ \ \ }{\tt\mc \ \ \ \ \ \ \ \ }\rm\mc {\tt /}{\tt *}\  {}の前にある。なぜ、そのように
					   しなければならない \hfill \rm\mc \ {\tt *}{\tt /}
\tt\mc 

\noindent
\mbox{}{\tt\mc \ \ \ \ \ \ \ \ }{\tt\mc \ \ \ \ \ \ \ \ }{\tt\mc \ \ \ \ \ \ \ \ }{\tt\mc \ \ \ \ \ \ \ \ }{\tt\mc \ \ \ \ \ \ \ \ }\rm\mc {\tt /}{\tt *}\  {}のか説明せよ。\hfill \rm\mc \ {\tt *}{\tt /}
\tt\mc 

\noindent
\mbox{}{\tt\mc \ \ \ \ }{\tt *}ptr[i]{\tt\mc \ }={\tt\mc \ }0;

\noindent
\mbox{}{\tt\char'175}

\noindent
\mbox{}\hfill

\noindent
\mbox{}void{\tt\mc \ }move{\tt\_}discs(int{\tt\mc \ }n,{\tt\mc \ }int{\tt\mc \ }i,{\tt\mc \ }int{\tt\mc \ }j,{\tt\mc \ }int{\tt\mc \ }k)

\noindent
\mbox{}{\tt\mc \ \ \ \ \ \ \ \ }{\tt\mc \ \ \ \ \ \ \ \ }{\tt\mc \ \ \ \ \ \ \ \ }{\tt\mc \ \ \ \ \ \ \ \ }{\tt\mc \ \ \ \ \ \ \ \ }\rm\mc {\tt /}{\tt *}\  {}上から $n$ 枚目までの disc
					   を、pole $i$ から pole $j$ に
					   \hfill \rm\mc \ {\tt *}{\tt /}
\tt\mc 

\noindent
\mbox{}{\tt\mc \ \ \ \ \ \ \ \ }{\tt\mc \ \ \ \ \ \ \ \ }{\tt\mc \ \ \ \ \ \ \ \ }{\tt\mc \ \ \ \ \ \ \ \ }{\tt\mc \ \ \ \ \ \ \ \ }\rm\mc {\tt /}{\tt *}\  {}pole $k$ を経由して、移動する
					   \hfill \rm\mc \ {\tt *}{\tt /}
\tt\mc 

\noindent
\mbox{}{\tt\char'173}

\noindent
\mbox{}{\tt\mc \ \ \ \ }if{\tt\mc \ }(n{\tt\mc \ }{\tt >}={\tt\mc \ }1)

\noindent
\mbox{}{\tt\mc \ \ \ \ \ \ }{\tt\char'173}

\noindent
\mbox{}{\tt\mc \ \ \ \ \ \ \ \ }{\tt\mc \ \ }move{\tt\_}discs(n{\tt\mc \ }{\tt -}{\tt\mc \ }1,{\tt\mc \ }i,{\tt\mc \ }k,{\tt\mc \ }j);{\tt\mc \ \ \ }\rm\mc {\tt /}{\tt *}\  {}関数 {\tt move\_discs()}
					   の中で、さらに自分自身 \hfill \rm\mc \ {\tt *}{\tt /}
\tt\mc 

\noindent
\mbox{}{\tt\mc \ \ \ \ \ \ \ \ }{\tt\mc \ \ }move{\tt\_}one{\tt\_}disc(i,{\tt\mc \ }j);{\tt\mc \ \ }{\tt\mc \ \ \ \ \ \ \ \ }\rm\mc {\tt /}{\tt *}\  {}{\tt move\_discs()} が使われ
					   ている。このような \hfill \rm\mc \ {\tt *}{\tt /}
\tt\mc 

\noindent
\mbox{}{\tt\mc \ \ \ \ \ \ \ \ }{\tt\mc \ \ }print{\tt\_}result();{\tt\mc \ \ \ \ \ \ \ }{\tt\mc \ \ \ \ \ \ \ \ }\rm\mc {\tt /}{\tt *}\  {}手法は、「再帰的呼びだし」
					   といわれる。一見 \hfill \rm\mc \ {\tt *}{\tt /}
\tt\mc 

\noindent
\mbox{}{\tt\mc \ \ \ \ \ \ \ \ }{\tt\mc \ \ }move{\tt\_}discs(n{\tt\mc \ }{\tt -}{\tt\mc \ }1,{\tt\mc \ }k,{\tt\mc \ }j,{\tt\mc \ }i);{\tt\mc \ \ \ }\rm\mc {\tt /}{\tt *}\  {}複雑に見える問題でも、
					   再帰的手法を用いると、\hfill \rm\mc \ {\tt *}{\tt /}
\tt\mc 

\noindent
\mbox{}{\tt\mc \ \ \ \ \ \ }{\tt\char'175}{\tt\mc \ }{\tt\mc \ \ \ \ \ \ \ \ }{\tt\mc \ \ \ \ \ \ \ \ }{\tt\mc \ \ \ \ \ \ \ \ }\rm\mc {\tt /}{\tt *}\  {}簡単に解けてしまうことが、
					   しばしばある。\hfill \rm\mc \ {\tt *}{\tt /}
\tt\mc 

\noindent
\mbox{}{\tt\char'175}

\noindent
\mbox{}\rm\mc {\tt /}{\tt *}\  \begin{center}
\includegraphics[scale=0.3]{hanoi1}\quad \includegraphics[scale=0.3]{hanoi2}
\end{center}


\noindent
たとえば、関数 {\tt move\_discs(4,i,j,k)} を呼び出すことは、
上図のような操作をすることに対応する。\hfill \rm\mc \ {\tt *}{\tt /}
\tt\mc 

\noindent
\mbox{}\hfill

\noindent
\mbox{}int{\tt\mc \ }main(void)

\noindent
\mbox{}{\tt\char'173}

\noindent
\mbox{}{\tt\mc \ \ \ \ }ptr[0]{\tt\mc \ }={\tt\mc \ }disc[0]{\tt\mc \ }+{\tt\mc \ }ARRAY;

\noindent
\mbox{}{\tt\mc \ \ \ \ }ptr[1]{\tt\mc \ }={\tt\mc \ }disc[1];

\noindent
\mbox{}{\tt\mc \ \ \ \ }ptr[2]{\tt\mc \ }={\tt\mc \ }disc[2];

\noindent
\mbox{}\hfill

\noindent
\mbox{}{\tt\mc \ \ \ \ }init{\tt\_}array();

\noindent
\mbox{}{\tt\mc \ \ \ \ }move{\tt\_}discs(ARRAY,{\tt\mc \ }0,{\tt\mc \ }1,{\tt\mc \ }2);{\tt\mc \ }{\tt\mc \ \ \ \ \ \ \ \ }\rm\mc {\tt /}{\tt *}\  {}{\tt ARRAY} 枚の disc を
					   pole 0 から pole 1 に pole 2
					   \hfill \rm\mc \ {\tt *}{\tt /}
\tt\mc 

\noindent
\mbox{}{\tt\mc \ \ \ \ \ \ \ \ }{\tt\mc \ \ \ \ \ \ \ \ }{\tt\mc \ \ \ \ \ \ \ \ }{\tt\mc \ \ \ \ \ \ \ \ }{\tt\mc \ \ \ \ \ \ \ \ }\rm\mc {\tt /}{\tt *}\  {}を経由して、移動
					   する \hfill \rm\mc \ {\tt *}{\tt /}
\tt\mc 

\noindent
\mbox{}{\tt\mc \ \ \ \ }return{\tt\mc \ }0;

\noindent
\mbox{}{\tt\char'175}

\noindent
\mbox{}

\rm\mc

\end{document}
